\section{Integrating a conditional sampling error}

Let $A$ denote the earlier successful histories and $E$ the actual
failure relation between a history and a fresh sampling variable.
A measurable relation can be integrated directly; a measurable choice
of auxiliary success sets for varying histories is unnecessary.

\begin{lemma}[Uniform section bounds]
    \label{lem:guarded_sampling_sections}
    \lean{MeasureTheory.Measure.prod_inter_fst_preimage_le_of_sections,
        MeasureTheory.Measure.prod_union_fst_compl_le_of_sections,
        MeasureTheory.Measure.prod_union_fst_compl_le_add_of_sections}
    \leanok
    Let $\mu$ be a measure on $X$ and $\nu$ an s-finite measure on $Y$.
    Let $A\subseteq X$ and $E\subseteq X\times Y$ be measurable.
    Write $E_x=\{y:(x,y)\in E\}$ and suppose that
    $\nu(E_x)\le\varepsilon$ for every $x\in A$, where
    $\varepsilon\in[0,\infty]$. Then
    \[
      (\mu\otimes\nu)(E\cap\pi_X^{-1}(A))
         \le\varepsilon\mu(A).
    \]
    If $\nu$ is a probability measure, then
    \[
      (\mu\otimes\nu)(\pi_X^{-1}(A^{\mathrm c})\cup E)
         \le\mu(A^{\mathrm c})+\varepsilon\mu(A).
    \]
    If $\mu$ is also a probability measure, the right side is at most
    $\mu(A^{\mathrm c})+\varepsilon$.
\end{lemma}
\begin{proof}\leanok
    By Tonelli's theorem, the first measure is the integral of the
    measures of its sections. Its integrand is bounded by
    $\varepsilon\mathbf1_A$. For the second assertion, split the
    failure set into the earlier-failure cylinder and the new failure
    over $A$. The cylinder has measure $\mu(A^{\mathrm c})$ and the
    first assertion bounds the remaining part. Finally $\mu(A)\le1$
    for a probability measure.
\end{proof}

These inequalities apply to the fixed-input estimates of
\cite[Lemma \texttt{l:rs}, lines 530--558, and merge analysis,
lines 1777--1817]{Arad2016RigorousRG}. Application to a complete
randomized interval tree additionally requires the measurable actual
parent operation and its uniform section estimate on earlier
successful histories.


\begin{lemma}[Combining the errors at a binary merge]
    \label{lem:binary_sampling_sections}
    \lean{MeasureTheory.Measure.prod_compl_prod_le,
        MeasureTheory.Measure.prod_binary_failure_le_of_sections}
    \leanok
    Let $\mu$, $\nu$ and $\rho$ be probability measures on measurable
    spaces $X$, $Y$ and $Z$, respectively. Let $A\subseteq X$ and
    $B\subseteq Y$ be the measurable successful child histories, and let
    $E\subseteq(X\times Y)\times Z$ be the measurable parent failure
    relation. Suppose that $\rho(E_{(x,y)})\le\varepsilon$ whenever
    $(x,y)\in A\times B$. Then
    \begin{align}
      ((\mu\otimes\nu)\otimes\rho)
       \bigl(\pi_{X\times Y}^{-1}((A\times B)^{\mathrm c})\cup E\bigr)
      \le \mu(A^{\mathrm c})+\nu(B^{\mathrm c})+\varepsilon.
    \end{align}
    Without any measurability assumption on $A$ and $B$, their outer
    measures satisfy
    \begin{align}
      (\mu\otimes\nu)((A\times B)^{\mathrm c})
        \le\mu(A^{\mathrm c})+\nu(B^{\mathrm c}).
    \end{align}
\end{lemma}
\begin{proof}\leanok
    \uses{lem:guarded_sampling_sections}
    The complement of $A\times B$ is the union of
    $A^{\mathrm c}\times Y$ and $X\times B^{\mathrm c}$.
    Subadditivity gives the second assertion. Apply the uniform section
    bound to the successful set $A\times B$ to obtain the first.
\end{proof}


\begin{lemma}[A uniform section bound in a finite product]
    \label{lem:finite_product_sampling_sections}
    \lean{MeasureTheory.Measure.pi_le_of_coordinate_sections}
    \leanok
    Let $\mu_j$ be probability measures on measurable spaces $X_j$ for
    $0\le j\le n$, and let $E\subseteq\prod_{j=0}^n X_j$ be measurable.
    Fix a coordinate $i$. For $x_{-i}\in\prod_{j\ne i}X_j$, let
    $E_{x_{-i}}$ be the set of values of the $i$th coordinate that complete
    $x_{-i}$ to an element of $E$. If
    $\mu_i(E_{x_{-i}})\le\varepsilon$ for every $x_{-i}$, then
    $(\bigotimes_{j=0}^n\mu_j)(E)\le\varepsilon$.
\end{lemma}
\begin{proof}\leanok
    \uses{lem:guarded_sampling_sections}
    Separate the $i$th coordinate from the remaining product. The
    canonical rearrangement preserves the product probability law.
    Apply the uniform section bound with all remaining-coordinate
    histories as the successful set, whose measure is one.
\end{proof}

\begin{theorem}[Uniform sections at any finite coordinate]
    \label{thm:finite_product_update_sections}
    \lean{MeasureTheory.Measure.pi_le_of_update_sections}
    \leanok
    Let $I$ be a finite index set and let $\mu_j$ be a probability
    measure on a measurable space $X_j$ for each $j\in I$.
    Fix $i\in I$ and a measurable set
    $E\subseteq\prod_{j\in I}X_j$. For a whole tuple $\omega$, write
    $\omega[i\leftarrow u]$ for the tuple obtained by replacing its
    $i$th coordinate by $u$. If every whole tuple satisfies
    \begin{align}
        \mu_i\{u:\omega[i\leftarrow u]\in E\}\le\varepsilon,
    \end{align}
    then $(\bigotimes_{j\in I}\mu_j)(E)\le\varepsilon$.
    No nonemptiness assumption on the $X_j$ is needed.
\end{theorem}
\begin{proof}\leanok
    \uses{lem:finite_product_sampling_sections}
    Since $i\in I$, enumerate $I$ by $\{0,\ldots,n\}$ for some $n$.
    This coordinate rearrangement preserves the product law. Every
    coordinate space is nonempty because it carries a probability
    measure. Thus a tuple of the remaining coordinates can be filled
    at coordinate $i$ to give a whole tuple $\omega$. Replacing that
    entry by $u$ gives precisely the section in the hypothesis.
    Apply Lemma~\ref{lem:finite_product_sampling_sections} and transport
    the resulting measure bound back through the enumeration.
\end{proof}

\begin{theorem}[A finite list of coordinate failures]
    \label{thm:finite_list_update_failures}
    \lean{MeasureTheory.Measure.pi_biUnion_list_le_of_update_sections}
    \leanok
    Let $\mu_j$ be probability measures on measurable spaces $X_j$,
    indexed by a finite set $I$. Let $l$ be a finite list of labels
    from an arbitrary set $K$, assign a coordinate $c(a)\in I$ to
    each label, and let $E_a\subseteq\prod_{j\in I}X_j$ be measurable
    for each label occurring in $l$. Suppose that each listed label
    and every whole tuple $\omega$ satisfy
    \begin{align}
        \mu_{c(a)}\{u:\omega[c(a)\leftarrow u]\in E_a\}\le\beta.
    \end{align}
    Then
    \begin{align}
        \left(\bigotimes_{j\in I}\mu_j\right)
          \left(\bigcup_{a\in l}E_a\right)\le |l|\beta.
    \end{align}
    Repeated labels are permitted and are counted in the list length
    $|l|$. The empty list is included.
\end{theorem}
\begin{proof}\leanok
    \uses{thm:finite_product_update_sections}
    Theorem~\ref{thm:finite_product_update_sections} bounds each listed
    event by $\beta$. Induction on the list and the inequality
    $\mu(A\cup B)\le\mu(A)+\mu(B)$ give the bound, counting every
    occurrence whether or not its event appeared earlier.
\end{proof}
